\documentclass[../main.tex]{subfiles}
\begin{document}
%==========================================TIÊU ĐỀ TẬP========================================
\begin{table}[h]
    \begin{adjustwidth}{-1cm}{}
        \begin{tabular}{>{\centering\arraybackslash}p{6cm}|>{\raggedright\arraybackslash}p{12.5cm}}
            \multirow{5}{6cm}
            % Cột bên trái: ÉPISODE và số tập
            {\rainbowcolor
                {\\[-40pt]
                    \flaregothic\fontsize{20pt}{20pt}\selectfont \centering
                    \color{eptype}JOUR\color{black}\\[12pt]
                    \hbrk\fontsize{72pt}{72pt}\selectfont
                    \color{epnum}8
                    \\[18pt]
                    % \flaregothic\fontsize{14pt}{14pt}\selectfont
                    % \color{epattr}BIENVENUE, \uline{\color{Banpire} Mel} !
                    % \flaregothic\fontsize{18pt}{18pt}\selectfont
                    % \color{epattr}ÉPISODE PERDU\\
                    \unrainbowcolor}
                \color{black}}
             &
            % Dòng thứ nhất: tên tập tiếng Nhật
            {\fotskip\fontsize{18pt}{18pt}\selectfont \ruby{葛}{くず}\ruby{葉}{は}\ruby{先}{せん}\ruby{生}{せい}の\ruby{学}{がく}\ruby{力}{りょく}\ruby{テ}{て}\ruby{ス}{す}\ruby{ト}{と}
            } \\[6pt]
            % Dòng thứ hai: tên tập tiếng Anh
             & {\cafeteria\fontsize{18pt}{18pt}\selectfont No Cheating Allowed!}                                                                                \\[4pt]
            % Dòng thứ ba: tên tập tiếng Việt
             & {\vnmsans\fontsize{18pt}{18pt}\selectfont Bài đánh giá năng lực}                                                                                 \\[8pt]
            % Dòng thứ tư: Nhân vật xuất hiện
             & {\sen{\fontsize{14pt}{14pt}\selectfont Track used: Chemistry (from Professional Block)}}
            \\[6pt]
            % % Dòng cuối cùng: Thời gian diễn ra của tập (thường sẽ là ngày phát hành)
             & {\continuum\fontsize{16pt}{16pt}\selectfont Ngày phát hành: 11/09/2021}                                                                          \\[18pt]
        \end{tabular}
    \end{adjustwidth}
\end{table}
%=========================================TRUYỆN CHÍNH========================================
\color{black}

Ngày thứ tám ở vùng đất Cầu Vồng của nhà Hikawa và thực thể Game.

\begin{wrapfigure}[26]{r}{6.5cm}
    \includegraphics[width=6.5cm]{../../../../Image/Book?/kuzuha_model.png}
\end{wrapfigure}

Tiết trời trong xanh một cách lạ kỳ. Nhưng không khí trong lớp lại như\dots\ nặng trịch. Bạn không nghe nhầm đâu, là trong lớp học, không phải ở trong văn phòng hay bât cứ đâu ở bên ngoài như thường lệ.

Kuzuha đứng thẳng người trước bục giảng, hai tay đút túi quần đồng phục, mặt không cảm xúc như thường lệ. Thế rồi, cậu nở một nụ cười ma mãnh, nói với giọng khản đặc: ``Hôm nay chúng ta sẽ làm bài kiểm tra nha.''

Chỉ một câu nói nhẹ tênh như thế, nhưng âm vang của nó khiến không ít người giật mình.

\chrint

Kuzuha (\ruby{葛}{くず}\ruby{葉}{は}, \ruby{Cư-dư}{Cát}-\ruby{ha}{Diệp}) từng là một trong những thành viên thuộc nhánh GAMERS, cùng với Yuika và Ririmu. Cậu đã từng bắt đầu công việc VTuber với tư cách độc lập, trước khi chính thức gia nhập Nijisanji vào tháng Bảy năm 2018 bốn tháng sau đó. Cậu cũng là Liver thành công nhất trong gia đình Nijisanji, và cũng là một trong những streamer nam tốt nhất trong giới YouTuber ảo.

Cậu là một chàng ma cà rồng với hồ sơ khá thảm hại: một tên NEET (\footnote{Viết tắt của ``Not in Edcuation, Employment or Training'' tức là những người dạng\\3 ``không'': không học hành, không việc làm, không đào tạo.}), vẫn còn sống ăn bám từ bố mẹ vốn dư dả tài chính và lúc nào cũng cần đến tiền. Mặc dù ngoại hình của cậu trông cũng không đến mức ``xấu xí'' bởi cậu là một ma cà rồng, nhưng bên trong lại ``đổ nát'': trẻ con, ích kỷ và đôi khi còn làm quá mọi chuyện lên. Cậu được các thành viên nhà Cầu Vồng coi là một người rất ồn ã và hay tấu hài cho mọi người. Tính cách làm quá mọi chuyện còn được thể hiện ở việc khi cậu chơi điện tử, cậu hay la hét hoặc tỏ vẻ tự mãn với kỹ năng của mình.

Tuy nhiên, trái với tính cách có phần tiêu cực kia, cậu là một người có khả năng đồng cảm với mọi người, từ đó có thể dễ dàng thích nghi với tình huống, thể hiện ở một lần cậu từng dỗ Chihiro theo cách của mình khi cô ấy bật khóc trên sóng. Khi cậu làm chung với các thành viên khác, cậu có xu hướng dè dặt hơn, dù sau này ``mọi chuyện lại đâu vào đấy'' khi cậu dần quen với đối phương hơn. Cậu cũng là một người đôi lúc sẽ dành lời khuyên cuộc sống cho người xem của mình. Cậu thậm chí đôi lúc còn nghi hoặc về khả năng của mình và thường sẽ làm cách nào đó để cải thiện bản thân hơn.

Dẫu cho cậu là một ma cà rồng, những trò chơi kinh dị lại\dots\ không phải là sở trường của cậu, khi cậu dường như dễ bị kích động và thường tìm cách ẩn mình càng câu càng tốt nếu có thể. Đây cũng là điểm khiến cậu bị mọi người trong công ty và cả những người hâm mộ thường xuyên dựa vào để châm chọc.

Kuzuha có mái tóc bạc trắng, đôi mắt đỏ mà ta thường bắt gặp ở những con ma cà rồng, dưới mắt trái cậu có một nốt ruồi, hai tai cậu đeo khuyên vòng nhỏ. Cậu mặc một chiếc áo khoác đen kín tới cổ, một chiếc quần dài màu xám có dây buộc phía trước, và thường đi dép lê trong nhà màu đỏ. Nói cách khác, một bộ trang phục mà những con lười hay diện khi ở nhà.

\chrint

Ngay khi Kuzuha dứt lời, đó cũng chính lúc mà cả lớp, bao gồm bốn gương mặt khác nhà Cầu Vồng và cặp đôi nhà Hikawa - Kuon và Kaigou - phản ứng với thông báo ấy theo những cảm xúc\dots\ rất khác nhau.

\img{nj-8a}

Về phần cậu con trai, cậu vỗ tay lách cách một cách chậm rãi, ánh mắt đầy mỉa mai: ``Yay\~{} Bài kiểm tra\~{} Vui quá xá là vui\~{}'' còn cô nàng thì thở dài thườn thượt, đầu úp xuống bàn: ``Mình ghét bài kiểm tra\dots''

Còn vì sao họ lại có mặt ở lớp thay vì ngồi trong công ty lập trình như thường lệ?

Chuyện là, từ tối hôm trước, sau khi Kuon, Kaigou và Suzuki với binh đoàn bốn người trở về đất liền, lúc đó khi hai người đang sắp dọn cơm tối sau một cuộc hành trình vất vả thì bất ngờ nhận được một tin nhắn từ Kuzuha: ``Ngày mai đến trường làm bài đánh giá năng lực.''

Không cần biết cậu ta biết được số của hai ``anh em'' nhà Hikawa, cũng không giải thích gì thêm. Chỉ vỏn vẹn một dòng tin như thế.

Cả Kuon lẫn Kaigou đều nhìn nhau với ánh mắt hoang mang xen lẫn bực bội.

``\vie{Hắn ta lại giở trò gì đây?}'' cô cằn nhằn. ``\vie{Mà sao hắn biết được số của chúng ta nhỉ?}''

``\vie{Chắc là mấy đứa chim lợn ở Nijisanji cho hắn số, rồi rảnh quá nên muốn kiếm trò gì đó giết thời gian đấy mà,}'' cậu đáp, giọng đầy mỉa mai.

Mặc dù cả hai không thoải mái gì với lời mời bất thình lình đó, họ vẫn chấp nhận nó. Có lẽ một phần là vì tò mò không rõ rốt cuộc tay ma dơi kia đang tính bày trò gì. Cũng có thể\dots\ vì thực sự họ cũng chẳng có gì để làm vào sáng hôm sau.

Theo tin nhắn của Kuzuha gửi cho họ sau đó, cậu gọi tổng cộng khoảng hai mươi người tới tham gia, gồm cả bốn thành viên khác thuộc nhà Cầu Vồng và hơn chục người thường mà cậu bắt được ở ngoài.

Suzuki thì được miễn. Kuzuha nói rằng ``bài kiểm tra này có phần hơi quá sức với cô bé'', nên cô xung phong ở lại ký túc xá để dọn dẹp, một việc cô bé cảm thấy phù hợp hơn.

Phản ứng của cặp đôi Hikawa này là thế, vậy các thành viên khác thì sao? Câu trả lời nằm ở bốn bức ảnh phía dưới.

\begin{figure}[H]
    \centering
    \begin{subfigure}[t]{0.45\textwidth}
        \centering
        \includegraphics[width=8cm]{../../../../Image/Book?/nj-8b1.png}
        \caption{Với Lize - một người có trình tri thức cao, đây là thông tin vừa bất ngờ vừa thú vị.}
    \end{subfigure}
    \quad
    \begin{subfigure}[t]{0.45\textwidth}
        \centering
        \includegraphics[width=8cm]{../../../../Image/Book?/nj-8b2.png}
        \caption{Với Mashiro, cậu chàng này hét lên thất thanh: ``\textbf{Bài kiểm tra bất ngờ ư!?}'', nhưng không tỏ vẻ sợ hãi gì.}
    \end{subfigure}
    \quad
    \begin{subfigure}[t]{0.45\textwidth}
        \centering
        \includegraphics[width=8cm]{../../../../Image/Book?/nj-8b3.png}
        \caption{Joe thì toát đầy mồ hôi, mắt trợn ra, cổ họng nghẹn lại không thốt lên lời.}
    \end{subfigure}
    \quad
    \begin{subfigure}[t]{0.45\textwidth}
        \centering
        \includegraphics[width=8cm]{../../../../Image/Book?/nj-8b4.png}
        \caption{Shouichi thì chẳng phản ứng, tay chống má, mắt vẫn nhắm như chưa hề đả động gì tới mình.}
    \end{subfigure}
\end{figure}

Tiện thể tới đây, ta cũng sẽ giới thiệu từng thành viên tham gia vào ``cuộc thí nghiệm xã hội'' - theo lời của Kuon và Kaigou.

\chrint

\begin{wrapfigure}[28]{l}{8.5cm}
    \includegraphics[width=8.5cm]{../../../../Image/Book?/lize_model.jpg}
\end{wrapfigure}

Đầu tiên là Lize Helesta (リゼ・ヘルエスタ, Li-dê He-lét-xta), gia nhập Nijisanji vào quý 1 năm 2019. Cô là công chúa thứ của Vương quốc Helesta, cũng là một người có trình tri thức cao. Cô giữ chức hiệu trưởng ở một trường học, một người đánh giá là ``văn võ vẹn toàn'', và cũng là một người siêng năng, tốt bụng với mọi người xung quanh, nên cô được mọi người yêu mến.

Theo mô tả của mọi người, Lize là một ``công chúa nhiệt huyết, xuất sắc cả về văn chương lẫn quân sự'', một người rất coi trọng việc rèn luyện hàng ngày và giao tiếp với mọi người để kế vị ngai vàng. Cô cũng trân trọng mối quan hệ với mọi người và dễ dàng rơi nước mắt, đặc biệt là khi phải chia tay. Cô duy trì hình ảnh nghiêm trang và đúng mực, nhưng đôi khi lại trở nên vui tươi hơn, đặc biệt là trong các trận đấu tập hoặc khi chơi game đối kháng, khi mà đôi lúc cô cũng sẽ bí mật chửi thề và dễ nổi cáu khi chơi game.

Mặc dù là một công chúa, cô lại có tính cách rất ngay thẳng, ghét sự xa hoa và luôn cố gắng tiết kiệm trong mọi việc. Cô thường chỉ ra sự lãng phí khi nhìn thấy, một điều không mấy phù hợp với hình ảnh của một công chúa. Tuy nhiên, cô lại không hề do dự khi nạp tiền vào các game gacha, đến nỗi em trai cô phải cảnh báo cô không được nạp quá nhiều.

Do địa vị công chúa của mình, Lize hầu như không có cơ hội kết bạn, khiến cô trở nên vụng về trong các mối quan hệ và có rất ít bạn bè. May mắn thay, cô có một người bạn từ thuở nhỏ là Anjou Katrina\footnote{Nhân vật này không xuất hiện trong bộ sách.}, một nhà giả kim thuật đến từ cùng ngọn núi. Khi còn nhỏ, Anjou thường xuyên đến vương quốc và chơi với cô, dần dần trở thành bạn thân. Chính vì thế mà dù chênh lệch tuổi tác tới chín năm, họ vẫn có thể nói chuyện như những người bạn cùng trang lứa.

Trên vai Lize luôn có một chú gà con là quản gia tên Sebastian Piyodore. Thực chất, đây là một con rối được điều khiển từ xa bởi máy móc chứ không phải sinh vật sống. Không ai rõ nó làm gì với tư cách một quản gia, nhưng mọi người đều xem nó như một linh vật đáng yêu của nàng công chúa này.

Lize diện bộ đồng phục của trường Nữ sinh Helesta, với  áo sơ mi trắng dài tay, một chiếc váy xếp ly màu xanh nhạt với váy ngoài cạp cao hai hàng khuy dưới ngực, đi quần tất màu xanh có dây đai đùi màu đen ở hai bên và bốt cao đến đầu gối buộc dây màu trắng. Cô có mái tóc trắng, nhuộm xanh lam ở bên trong và có một cái kẹp tóc đính ở bên trái, mang hình hai chiếc lá màu xanh đậm và xanh lam xếp chồng lên nhau.

\begin{wrapfigure}[30]{r}{11cm}
    \includegraphics[width=11cm]{../../../../Image/Book?/mashiro_model.jpg}
\end{wrapfigure}

Tiếp đến là người ngồi cạnh cô, Mashiro Meme (ましろ\ruby{爻}{めめ}, Ma-si-rô \ruby{Mê-mê}{Hào}), một cậu chàng (xin nhấn mạnh, là cậu chàng, mặc dù ngoại hình rất dễ bị hiểu nhầm) gia nhập vào công ty Cầu Vồng vào quý 4 năm 2019. Cậu mang đậm hình ảnh một kẻ điên khùng, khi lúc nào cũng cầm trong tay một chiếc xẻng. Khi một khán giả hỏi cậu dùng xẻng để làm gì, cậu chỉ chưng hửng trả lời: ``Để chôn người'' và ``dùng nó ở nhiều nơi ngoài việc đào hố.'' Người ta thường nói rằng Mashiro có một tuổi thơ không mấy êm đẹp, và đích thân cậu ta đã nói thế.

Cậu chàng này là một sinh viên đại học nhanh nhẹn và tham gia vào mọi việc, và cũng là một người ưa mạo hiểm, không ngần ngại thử bất cứ điều gì mình thích. Sở thích của cậu bao gồm có kinh dị, những câu chuyện huyền bí và những hiện tượng siêu nhiên. Đôi khi, cậu còn thực hiện các nghi lễ ``trốn tìm một mình'' trên sóng trực tiếp, khiến nhiều người phải đổ mồ hôi hột. Chính sở thích đó nên cậu không tỏ ra nao núng khi chơi các tựa game kinh dị.

U ám và buồn bã cũng là những chủ đề mà cậu chàng này ưa thích. Mashiro thừa nhận rằng lý vì sao cậu bị cuốn hút và không hề sợ hãi những điều kinh dị hay siêu nhiên là vì anh ấy không tin vào sự tồn tại của chúng, và do đó chủ động tìm kiếm hoặc thách thức chúng để thỏa mãn trí tò mò của mình.

Cậu tự mô tả mình là người hướng nội và kém trong việc tiếp cận người khác, dẫn đến số lượng video hợp tác với các vTuber khác tương đối thấp. Mashiro cũng là một người có giọng bình thường nghe rất dịu dàng, nhưng khi hát lại rất khàn, hợp với những bài rock hơn là ballad.

Mashiro diện áo khoác màu xanh đậm, bên trong là áo sơ-mi trắng in hình trái tim (bị chảy máu?) bên ngực trái, đeo cà vạt xanh, quàng một chiếc khăn đỏ viền hoa quanh cổ, diện với chiếc quần đen có dây đeo, quần tất màu vàng có các chấm đen dọc xuống và dép xăng-đan màu đen. Gu ăn mặc của cậu chàng này khá là dị hợm, thậm chí còn vô tình ``khiến cảnh sát trang phục phải thấy ngứa mặt'' (lời bình của Kuon).

\begin{wrapfigure}[29]{l}{7cm}
    \includegraphics[width=7cm]{../../../../Image/Book?/joe_model.jpg}
\end{wrapfigure}

Ngồi bên cạnh cậu chàng đam mê với những thứ kinh dị là ``chú hề'' của Nijisanji, Joe Rikiichi (ジョー・\ruby{力}{りき}\ruby{一}{いち}, Giâu \ruby{Ri-ki}{Lực}-\ruby{i-chi}{Nhất}). Cậu vốn dĩ thuộc nhánh SEEDs, cùng với Ritsuki và Debiru\footnote{Sẽ xuất hiện ở {\flaregothic JOUR \hbrk 12}.} Tên của cậu hề này đọc khá giống với ``{\sen JOKER} (ジョーカー)'', bản thân cậu chàng cũng trông giống một chú hề trong ``Smile PreCure!''

Joe là một chú hề và cũng là một diễn viên hài, về cơ bản cậu ta là một người có chuyên môn là bắt chước giọng nói và ca hát, những điều mà cậu làm hơn hẳn so với những chú hề bình thường. Cậu cũng là người giỏi thích nghi và cách đối phó với các tình huống bất ngờ xuất hiện, và ngay cả khi các thành viên khác mắc lỗi hoặc phạm sai lầm, Joe có thể chấp nhận mọi thứ và biến chúng thành một trò đùa (mà mọi người có thể cười và không thấy người bị mắc lỗi cảm thấy bị xúc phạm). Dù là một chú hề, thu nhập của cậu cũng không được cao, điều dễ thấy trong các rạp xiếc.

Joe cũng có một kỹ năng ca hát đáng kể, với giọng hát có phổ rộn gvaf có thể hát trong khi hòa phối với các thành viên khác hoặc hát những nốt cao mà không gặp vấn đề gì.

Cậu là bạn thân của bạn cùng lớp Keisuke Maimoto và là người dẫn chương trình phát thanh đêm khuya. Ngoài ra, cậu cũng có thành lập một nhóm nhạc riêng trong Nijisanji có tên là ``Rain Drops'' (レインドロップス) cùng với Akina, Ryushen và ba thành viên khác. Họ thậm chí còn có chương trình radio và đã ra mắt ba album cho Universal Music.

Joe là một trong số ít những thành viên có phiên bản khác giới của mình, và ở đây là Rokumeikan Kiriko (\ruby{鹿}{ろく}\ruby{鳴}{めい}\ruby{館}{かん}キリコ, \ruby{Rô-cư}{Lộc}-\ruby{mây}{Minh}-\ruby{can}{Quán}\ Ki-ri-cô), một bản tính thay thế của chính cậu. Về cơ bản cô nàng này được xây dựng như mọt nữ phú bà giàu có, mũi cao, và có ngoại hình hao hao gống với Hyakumantenbara Salome\footnote{Nhân vật này không xuất hiện trong bộ sách.}, và tréo ngoe hơn, cả hai người này lại\dots\ khá thân thiết.

Joe thường trang điểm đậm, trong đó cậu đánh son môi màu xanh lục, phía mắt trái có vẽ một giọt nước màu hồng tím, còn mái tóc cậu thì trông khá bù xù và mang hai màu giống với màu môi và màu giọt nước: xanh lục và hồng tím. Bộ trang phục của cậu gồm một chiếc áo sơ-mi hồng, mặc bên trong bộ vest trắng, cà vạt tím in hình họa tiết con mắt, đeo thắt lưng đen, găng tay đen, diện với quần trắng, đi tất tím in hình mắt và giày da trắng.

\begin{wrapfigure}[28]{r}{9.5cm}
    \includegraphics[width=9.5cm]{../../../../Image/Book?/shoichi_model.png}
\end{wrapfigure}

Người cuối cùng, ngồi ở bên cửa ra vào là Kanda Shouichi (\ruby{神}{かん}\ruby{田}{だ}\ruby{笑}{しょう}\ruby{一}{いち}, \ruby{Can}{Thần}-\ruby{đa}{Điền}\ \ruby{Sô}{Tiếu}-\ruby{i-chi}{Nhất}), một thành viên thuộc nhánh SEEDs cùng với Keisuke và Rion\footnote{Sẽ xuất hiện ở {\flaregothic JOUR \hbrk 12}.}.

Một cậu chàng sinh viên 21 tuổi chuyên ngành quan hệ truyền thông. Cậu ta là một người thích trò chuyện với mọi người, nhưng đôi khi lại nói những lời cay nghiệt với người khác, chỉ để rồi sau đó cảm thấy hối hận với những gì cậu nói khi ``mình đã đi quá xa!''. Bất chấp như thế, Shouichi là một chàng trai có trái tim chân thành. Cậu dường như đang giữ một ranh giới mong manh giữa lý trí và tình cảm, dù rằng lại hay lo những việc khá vặt vãnh.

Với nền tảng từng học phát thanh, cậu có khả năng phát âm cực kỳ rõ ràng, được đánh giá là một trong những người có kỹ năng tường thuật tốt nhất tại Nijisanji. Không chỉ dừng lại ở đó, Shouichi còn học thêm về giải trí, và qua việc tham gia nhiều dự án voice drama, cậu đã chứng tỏ được khả năng diễn xuất cùng kỹ năng kể chuyện thuộc hàng top của công ty.

Một điều thú vị khác về Shouichi là cậu có khả năng đọc bài tarot. Có một lần, cậu từng bói cho một thính giả đã bị sa thải ở công việc trước đó. Khi rút được lá The Chariot (Con Chiến Xa), cậu đã nửa đùa nửa thật động viên: ``Sao không thử đến Ichikara?''. Điều bất ngờ là người nghe đó đã thực sự nộp đơn xin việc vào Ichikara ngay sau khi được bói bài và đã được nhận. Sự việc này khiến cả Shouichi lẫn các thính giả đều ngạc nhiên trước quyết tâm và động lực của người ấy.

Nhìn bề ngoài, Shouichi thường để đôi mắt nhắm, nhưng thực chất là mắt cậu mở hé, và hiếm khi mở to. Cậu có mái tóc vàng. Bộ trang phục của cậu bao gồm một chiếc áo trùm đầu màu xám mang họa tiết những hình thoi đỏ xếp đan chéo với các đường đen song song, một chiếc áo khoác đen bên ngoài với cổ tay có họa tiết ca-rô xanh lá, một chiếc quần capri đỏ (quần dài tới trước mắt cá chân) và giày thể thao màu xanh lá.

\chrint

``\vie{\hl{\dots Trông mọi người có vẻ lo lắng nhỉ,}}'' Kuon vừa nói vừa liếc nhìn quanh lớp, rồi lặng lẽ dựa người ra sau ghế, ánh mắt vẫn hướng về phía bục giảng nơi Kuzuha đang lúi húi lục ra một xấp giấy dày cộm.

``\vie{\hl{Cũng đúng thôi,}}'' cô nàng cười trừ, ánh mắt không rời khỏi thầy giáo bất đắc dĩ của lớp, ``\vie{\hl{tự dưng nghe kiểm tra mà không biết nội dung gì, ai mà chẳng hoang mang.}}

``\vie{\hl{Hầy\dots{} Đúng là gợi lại cảm giác cấp ba thật đấy,}}'' cậu khẽ thở ra một tiếng, ``\vie{\hl{mấy bài kiểm tra 15 phút, rồi 1 tiết, rồi cuối kỳ. Cứ tưởng đời sinh viên là hết rồi cơ.}}''

``\vie{\hl{Nhưng đúng là cái kiểu trò mà gã Kuzuha kia sẽ bày ra thật,}}'' cô chép miệng, ``\vie{\hl{gọi cả đám tụ tập lại, đùng một cái bảo kiểm tra đột xuất, ném cho cái đề, rồi để mặc bọn mình lo.}}''

``\vie{\hl{Cơ mà tôi cũng không thấy hắn nhắc tới điểm số, hình phạt hay phần thưởng gì cả.}}'' cậu khoanh tay, nhíu mày, ``\vie{\hl{chắc là đơn thuần để đánh giá năng lực từng người thôi.}}''

``\vie{\hl{Vậy thì có gì phải áp lực đúng không?}}'' cô nàng nhún vai, rồi hạ thấp giọng, ``\vie{\hl{nhưng tôi cũng muốn thử xem mình cỡ nào. Chẳng lẽ để bị bọn kia vượt mặt à?}}''

``\vie{\hl{Cô vẫn cái kiểu cạnh tranh ngầm nhỉ.}}''

``\vie{\hl{Gọi là động lực tinh thần thôi.}}''

``\vie{\hl{Tôi thì chẳng quan trọng mấy. Cứ làm cho xong là được.}}''

``\vie{\hl{Gớm\~{} Để xem lúc làm bài thì ai xong trước ai.}}''

``\vie{\hl{Tôi cũng chẳng muốn hơn thua đâu. Mệt lắm.}}'' cậu ngáp một cái, nhưng mắt vẫn nhìn chằm chằm tờ đề vừa được phát tới tay.

Và thế là, bài kiểm tra đánh giá năng lực của lớp Nijisanji lập vội chính thức bắt đầu.

Đó cũng là lúc những chiêu trò ``truyền thống'' như quay cóp, nhìn trộm, giả vờ đau bụng hay\dots\ tâm linh học đường bắt đầu được triển khai bởi những kẻ không học bài.

Nhưng liệu họ có thể lách luật trót lọt trước một giáo viên đầy mưu mẹo như Kuzuha không?

Câu trả lời sẽ đến rất sớm thôi.

\ct

Trong căn học rộng, hai mươi học sinh ngồi rải rác theo từng bàn đơn, với khoảng cách đủ xa để ngăn chặn mọi ý định quay cóp hay thì thầm. Gần một phần ba trong số đó là những người mà Kuon và Kaigou chưa từng gặp qua, trông có vẻ như dân thường hoặc nhân viên trong khu thành phố Cầu Vồng được mời đến để kiểm tra trình độ.

Toàn bộ không gian ngập trong thứ âm thanh quen thuộc mà bất kỳ lớp học nào cũng có thể sản sinh ra: tiếng giấy lật phạch phạch, tiếng bút chì gõ nhẹ vào mặt bàn hoặc xoèn xoẹt lia trên giấy, và vài tiếng thở khẽ nhè nhẹ, đủ để nếu ai chăm chú lắng nghe thì sẽ phân biệt được từng nhịp một và tiếng đó đến từ đâu.

Ngay khi lật tờ đề lên, đôi mắt của cặp anh em nhà Hikawa lập tức lia xuống những câu hỏi đầu tiên.

``\textit{Toán à\dots}'' cả hai cùng nghĩ thầm, cùng lúc, dù chẳng ai nói ra.

Một đề thi trắc nghiệm bốn lựa chọn, có yêu cầu viết lời giải tắt bên cạnh. Bốn mươi câu trong vòng một giờ, với bố cục chia đôi rõ ràng: hai mươi câu đại số và giải tích, cùng với 20 câu hình học phẳng và không gian.

``\textit{Dễ hơn cả đề cấp ba nữa\dots}'' Kuon nhận định thầm, gương mặt không biến sắc. Những câu đầu là các bài toán số học dạng tìm đạo hàm, một vài công thức lượng giác, chuỗi số và phương trình lũy thừa. Đây là vùng đất sở trường của cậu, người từng học giải tích bậc đại học như cơm bữa.

Trong khi đó, Kaigou khẽ mỉm cười khi thấy phần ba hình đầu tiên là tổ hợp các khối đa diện, mặt phẳng giao nhau, và phép chiếu. ``\textit{Đúng bài mình thích\dots}'' cô nàng nghĩ thầm, bút chì đã xoay nhẹ trong tay. Dù phần đại số không phải điểm mạnh, nhưng nếu là không gian ba chiều và những hình phẳng ảo ma tới độ người thường nhìn không hiểu, thì Kaigou đủ tự tin mình sẽ vượt trội hơn vài người trong lớp.

Cả hai người họ ngồi cùng bàn ở hàng đầu, ngay trung tâm lớp học, đối diện trực tiếp với bàn giáo viên, và khuôn mặt họ vẫn giữ khuôn mặt tập trung, nghiêm túc. Họ chọn vị trí đó không phải vì muốn chứng tỏ điều gì, mà đơn giản vì\dots\ ở đó có quạt trần mát, và họ vốn chẳng quan tâm người đứng lớp là ai. Với họ, ``gian lận'' là điều không cần thiết, và cuộc thi này dù bất ngờ cũng chẳng đe dọa gì đến danh tiếng hay công việc của cả hai.

Không liếc nhau một lần, hai bộ bút chì bắt đầu trượt đều trên mặt giấy trắng. Trong đầu mỗi người đều vang lên tiếng tích tắc vô hình, như những cuộc chạy đua song song nhưng không ganh đua.

Một bài kiểm tra tình cờ. Một lần trở về những ngày học sinh. Một cách để xác nhận trình độ.

Và với Kuon cùng Kaigou, đó là một dịp hiếm hoi để thử xem, bản thân họ đã đi xa tới đâu.

Vậy bốn thành viên còn lại của nhà Cầu Vồng thì sao?

Mashiro ngồi ở bàn cuối cùng, cùng hàng với Lize (bên trái) và Joe (bên phải). Shouichi thì ngồi ở bàn trước họ, ngay sát cửa ra vào. Cả ba đều đang ở trong tâm trạng\dots\ bất ổn, mỗi người một kiểu.

Ngay khi nhận đề kiểm tra, cậu bạn Kinh dị khẽ co rúm lại. ``\textit{Toán á!? Chết tôi rồi\dots}'' Chân tay cậu run lẩy bẩy, đôi mắt đảo quanh phi logic như đang tìm lối thoát khẩn cấp giữa một mê cung đề thi. Cậu gắng nhớ lại chút kiến thức từng học thời phổ thông, nhưng tất cả hiện lên chỉ là những vệt mờ mịt như sương sớm.

\img{nj-8c}

Sau vài phút cắn bút và toát mồ hôi, cậu nghĩ thầm:
``\textit{Không được rồi\dots\ Hay là mình thử liếc sang bên cạnh xem sao?}''

Ngay lập tức, ánh mắt cậu hướng về phía trái, nơi Lize đang miệt mài giải bài. Cô công chúa tóc bạch kim ấy dường như không màng đến xung quanh, đôi mắt sắc lạnh và tập trung, cây bút chì nhịp nhàng trượt theo từng dòng đề bài, không hề ngập ngừng.

\img{nj-8d}

Mashiro há hốc miệng: ``\textit{Quái thật!? Cô ấy làm như thể đây là bài kiểm tra tiểu học vậy\dots\ Mình biết Lize-sama học giỏi rồi, nhưng cỡ này thì\dots}''

``\textit{\dots Đành liều vậy.}'' Cậu rướn người ra hẳn bên trái, thân hình dài một cách bất tự nhiên, mắt đảo sang bài của Lize, căng não đọc thật nhanh những gì cô viết.

Và thế là, tay cậu Mạo hiểm bắt đầu cắm cúi chép lại từng dòng, từng đáp án, vừa chép vừa rón rén thở ra một tiếng nhẹ: ``\textit{Hú\~{}\dots Quá êm.}''

Joe ở bàn bên cạnh, thấy thế thì cũng làm theo. Dù chẳng hiểu gì mấy, gã vẫn cười hí hửng, liếc ngang bài của Mashiro, rồi bắt đầu ``copy'' y chang, chỉ khác là tay gã vung vẩy như múa, trông cứ như đang biểu diễn màn xiếc nho nhỏ trong giờ thi.

Ở phía trước, Shouichi thì dùng cách thô bạo hơn: rút điện thoại ra, lén chụp bài của gã Hề sau khi đã ``qua xử lý'' từ Mashiro, rồi tự mình chép lại. Thậm chí cậu còn ngáp nhẹ một cái, như thể bài kiểm tra này là tiết học Lịch sử buồn ngủ chứ không phải là phép thử năng lực.

\begin{figure}[H]
    \centering
    \begin{subfigure}[t]{0.45\textwidth}
        \centering
        \animategraphics[autoplay,loop,width=8cm]{30}{../../../../Image/Book?/nj-1c/nj-1c.}{1}{87}
    \end{subfigure}
    \quad
    \begin{subfigure}[t]{0.45\textwidth}
        \centering
        \animategraphics[autoplay,loop,width=8cm]{30}{../../../../Image/Book?/nj-1c/nj-1c.}{134}{234}
    \end{subfigure}
    \quad
    \begin{subfigure}[t]{0.45\textwidth}
        \centering
        \includegraphics[width=8cm]{../../../../Image/Book?/nj-8e3.png}
    \end{subfigure}
    \caption{Quy trình sao chép bài của Mashiro, Joe và Shouichi.}
\end{figure}

Về phía cặp đôi Hikawa thì sao? Họ dường như không mảy may về những gì đang diễn ra sau lưng, về một ``cỗ máy copy-paste'' đang hoạt động bí mật.

Cả hai người vẫn chăm chú làm bài, không nói một lời, không liếc nhìn nhau. Với cậu Tổng nhà hololive, các câu đại số - giải tích cứ thế trôi tuột qua tay. Dù vậy, phần hình học bắt đầu khiến cậu phải nhăn mặt đôi chút, nhất là mấy câu dựng hình hoặc tính góc giữa hai mặt phẳng.

``\textit{Tính góc giữa hai mặt\dots\ được rồi, cái này mình vẫn còn nhớ.}''

Ngược lại, Kaigou đang hì hục đè bút giải các khối đa diện, các mặt nón cắt bởi mặt phẳng xiên. Nhưng tới khi gặp một hệ phương trình hai ẩn, cô nàng cũng đã bắt đầu chùn bước đôi chút.

``\textit{Dùng định lý Vi-ét cơ bản\dots\ Được rồi, mình vẫn còn nhớ $\delta$ tính như thế nào. Vẫn trong tầm kiểm soát.}''

Ở bục giảng, Kuzuha đứng khoanh tay quan sát cả lớp. Bên ngoài trông cậu có vẻ thờ ơ, mắt lơ đãng, dáng lười nhác. Nhưng thực ra, cậu đang ghi nhớ từng cử động nhỏ nhất của các học sinh. Dù không cầm sổ theo dõi, đầu cậu đã lưu lại mọi thứ.

``\textit{Mashiro liếc ba lần, Joe nhìn sang trái năm lần, Shouichi vừa bấm nút điện thoại,}'' cậu thầm tổng kết trong đầu. ``\textit{Nhà Hikawa vẫn tập trung làm bài, bắt đầu chậm dần tốc độ. Hừm.}''

``\textit{Để xem ai đủ tinh vi, và ai đủ ngốc nào.}'' Nói là làm, Kuzuha lập tức rời khỏi bục giảng và bắt đầu đi tuần xung quanh mọi người, đi dần về phía bàn của Lize và Mashiro.

Khi cậu bạn Mạo hiểm liếc nhìn giám thị đang di chuyển khỏi hàng ghế cuối, cậu lập tức quay sang Lize, thân hình vươn dài lần nữa như một chú mèo xám trườn qua mặt bàn. Chép. Cúi đầu. Viết tiếp. Lặp lại.

Joe bắt chước lại lần nữa ngay khi cậu Ma cà rồng dần quay trở lại bàn giáo viên, hai tay lắc lư như để xem Mashiro đã bổ sung gì thêm.

Nhưng đúng lúc ấy\dots

*CỘP!*

Kuzuha không nói không rằng, bất chợt quay đầu lại và tiến tới Joe, và không thương tiếc đè đầu gã hề xuống bàn, khiến âm thanh vang lên trong không gian tĩnh mịch như tiếng búa gõ vào trống.

\img{nj-8g}

Gã tóc hai màu định ngẩng đầu lên, nhưng đã quá muộn.

Mashiro giật thót, nhưng cố ra vẻ như đang tra cứu công thức trong đầu, nhưng vẫn không khỏi cảm giác lạnh sống lưng khi bị giám thị bắt được tại chỗ.

Shouichi thì chống tay lên má, giả vờ suy nghĩ. ``\textit{Thằng đần này đúng là ngu thật}'' câhy thầm rủa, vì giờ đây cậu không còn ai để có thể lợi dụng được nữa.

Kuzuha chỉ thở dài, rồi lặng lẽ quay trở lại bục giảng. ``Joe-san bị cảnh cáo lần một.''

\ct

Không lâu sau đó, cặp đôi nhà Hikawa đã hoàn tất bài kiểm tra và soát lại bài đến BA lần. Mỗi lần kiểm tra là một lần rà soát lại cách trình bày, kiểm tra tô bút chì, rồi đối chiếu lại cách tư duy của chính mình. Từng câu từng chữ đều được rà tỉ mỉ.

Với Kuon, dù là người ưa số học, nhưng vài câu hình học vẫn là bài toán hóc búa. Có những câu buộc cậu phải khoanh bừa trong khuôn khổ logic, tức không phải kiểu đoán mò, mà là phân tích sơ đồ, nhớ lại một chút công thức từ cấp ba, gán vào vài giả thuyết và\dots\ đánh cược. Cậu là người dạng luôn tính nhẩm, nên cậu không bao giờ xin tờ nháp.

Trong khi đó, Kaigou, người nổi trội với phần hình học, thì xử lý cực kỳ thận trọng các bài tập dạng bảng biến thiên, đồ thị hàm số và nguyên hàm theo những gì cô còn nhớ. Cô nàng xin thêm tới ba tờ giấy nháp để tính đi tính lại các bước, vẽ sơ đồ trục toạ độ chính xác đến từng góc. Phần đại số tuy là điểm yếu, nhưng cô cũng không để bị ``lụt mất thuyền''.

Và khi tất cả đã xong, cả hai người úp mặt xuống bàn ngủ một cách sảng khoái. Một giấc ngủ như phần thưởng cho khoảng thời gian nghiêm túc, không gian lận, không ngó nghiêng ai, mà họ cũng chẳng cần. Họ không phải thiên tài toàn diện, nhưng là những người hiểu rõ năng lực mình tới đâu và sẵn sàng dùng tất cả để chứng minh điều đó.

Ở phía sau, nàng Công chúa xứ Helesta cũng vừa kết thúc bài làm của mình. Nhẹ nhàng, cô vẽ một chú chim non, biểu tượng của quê nhà, vào góc trên của tờ giấy, rồi đặt bút chì xuống, giãn vai thật sâu và nở một nụ cười tự hào.

\img{nj-8h}

Chính hành động đó lại trở thành điểm yếu chí mạng. Nhưng không phải dành cho cô.

Mashiro - vốn đang âm mưu thực hiện bước cuối của kế hoạch ``copy sạch mọi thứ'' - vừa trông thấy Lize không để ý, lập tức ngả người sang bàn bên, rút lấy bài làm của cô nàng, rồi hí hửng chép lại toàn bộ bài làm với tốc độ thần sầu. Chỉ trong vòng chưa đầy hai phút, cậu đã chép xong nửa bài còn lại, rồi nhanh chóng trả bài về chỗ cũ, ngay đúng lúc Lize vừa ngáp một cái thật dài và quay đầu lại.

\gif{nj-8i}{1}{131}

Hai bên nhìn nhau.

Cô Công chúa cười nhẹ vì cảm giác sảng khoái.

Cậu Mạo hiểm cũng cười nhẹ vì đã thực hiện được ý đồ của mình.

Cả hai không ai phát hiện ra điều gì. Hoặc ít nhất, họ nghĩ là vậy.

Ở bục giảng, Kuzuha liếc qua toàn bộ lớp một lần nữa, đôi mắt lạnh như thủy tinh quét khắp căn phòng. Sau đó, cậu lên tiếng: ``Hết giờ. Tất cả nộp bài lên bàn giáo viên.''

Mọi người bắt đầu đứng dậy. Nhưng hai người ở bàn đầu - Kuon và Kaigou - vẫn đang ngủ say như chết.

``\textit{Tôi nói cả hai cô cậu đấy, Hikawa-san!}'' chàng Ma cà rồng nâng giọng một chút.  Hai con người phía trước đồng loạt bật dậy như zombie tỉnh giấc, mắt còn chưa mở rõ.

``\vie{Ugh\dots\ Hết giờ rồi à\dots?}'' cậu rên rỉ, vươn hai tay lên cao.

``\vie{Nhanh thật đấy\dots}'' cô nàng lảo đảo, một tay chống cằm lên bàn, tay kia hướng sang bàn bên cạnh. ``\vie{Kuon-kun\dots\ Đưa bài cho tôi nào\dots}''

``\vie{Đây\dots}'' cậu Tổng đưa bài làm của mình, rồi ngáp một tiếng rõ to. ``\vie{Má, buồn ngủ quá\dots}''

``\vie{Mới chợp mắt được một tí thôi mà\dots}'' cô lấy bài từ Kuon, thở dài một hơi. Cô lững thững tới bàn giáo viên, nộp cả hai bài rồi quay về chỗ với vẻ mặt chẳng buồn che giấu sự mệt mỏi.

Ở phía cuối lớp, Mashiro vui sướng nâng bài làm của mình, đầu óc tràn ngập cảm giác chiến thắng. Trong khi đó, Joe thì mặt mày xanh lét, mồ hôi túa ra như vừa chạy bộ, tay run run cầm bài thi.

Shouichi vẫn chưa nộp, chống tay nhìn đời, cố giấu vẻ tuyệt vọng trong lòng. ``\textit{Xong rồi\dots\ Kiểu này là ăn trứng ngỗng thật rồi,}'' cậu thở dài, nhưng đúng lúc ấy, một con bướm bay ngang qua cửa sổ, khiến Lize đang mơ màng đưa mắt ra nhìn theo.

Và điều mà cậu thấy? Bài làm của cô nàng Công chúa, vẫn chưa được nộp.

Cậu lập tức đứng phắt dậy. ``Đây là cơ hội sống sót duy nhất!''

Trong tích tắc, Shouichi nhảy tới bàn của Lize, xóa tên trên bài cô, viết đè lên tên mình, rồi quay về bàn với tốc độ của một vận động viên điền kinh\dots\ nhưng hoàn toàn không để ý đến biểu tượng con chim nhỏ ở góc trên. Cậu cũng không quên ghi lại tên của Lize lên bài làm của mình.

\gif{nj-8j}{1}{154}

Joe chứng kiến mọi thứ như được gãi đúng chỗ ngứa, cũng nghĩ thầm: ``\textit{Ý hay đấy.}''

Gã nhanh chóng viết tên Shouichi lên bài của mình, rồi lấy bài của Lize, xóa dòng tên đè kia đi và viết tên mình vào, y như Shouichi\dots\ và y như một tên trộm hậu đậu không để ý đến dấu vết rõ ràng nhất: con chim non kia.

\gif{nj-8k}{1}{234}

Nàng Công chúa khi đứng dậy nộp bài, cô cầm lấy tờ giấy gần nhất mà cô cho là của mình, không để tâm gì nhiều. Mashiro, Shouichi và Joe cảm thấy mình vừa hoàn thành một ``phi vụ tráo bài'' thần thánh. Cậu bạn Truyền thông thậm chí còn vừa cười vừa giương ánh mắt đắc chí về phía trước, một điều mà cậu hiếm khi làm.

Cả lớp dần nộp xong bài. Những tiếng thì thầm lắng xuống. Mọi hành động mờ ám - ít nhất là trong tâm trí họ - đều trót lọt.

Nhưng Kuzuha thì sao?

Vẫn đứng yên sau bàn giáo viên. Đôi mắt xám ấy chỉ khẽ khép lại. ``Thú vị thật đấy.''

\ct

Giờ nghỉ giữa giờ. Tất cả mọi người có một khoảng thời gian xả hơi sau sáu mươi phút căng thẳng bài kiểm tra.

Kuzuha dựa lưng vào ghế giáo viên, chậm rãi lật từng trang bài kiểm tra của đám học sinh hôm nay. Cậu vốn dĩ đã để mắt tới ít nhất bốn người trong số đó - ba kẻ láu cá ngồi cuối lớp và một cô công chúa nghiêm túc - nên khi nhìn thấy bài thi ghi tên Mashiro, cậu chỉ khẽ nheo mắt.

``Ừm\dots\ nét chữ này thì đúng là của Mashiro thật.''

Bài làm rất chỉn chu, lời giải ngắn gọn, tốc độ làm bài nhanh đến đáng ngờ, đặc biệt với một thằng nhóc từng run rẩy khi cầm đề. Cậu lật lại trí nhớ, hình ảnh cậu bạn Mạo hiểm kéo dài người liếc sang bàn bên cạnh lại hiện ra trong đầu như một thước phim quay chậm.

``Cái trò biến thành dây thun ấy\dots\ đúng là hết thuốc chữa.''

Cậu đặt bài xuống, mở sang bài ghi tên là Joe - cũng chính là bài gốc của Lize. Cách trình bày đều đặn, câu nào cũng có lý giải rõ ràng, nét bút thẳng hàng, dễ đọc như sách in. Cậu khẽ nhướng mày khi nhìn thấy một số câu bị giải sai, hoặc là những câu bỏ trống phần giải thích.

``Còn cái này thì\dots\ là của Lize, khỏi nghi ngờ. Hình chim non kia là đủ hiểu rồi.''

Kuzuha cầm cây bút đỏ, khoanh tròn hình vẽ đó một cách tàn nhẫn, như thể tuyên án tử hình cho âm mưu dối trá của những kẻ không học bài. ``Joe à, ít nhất cũng học cách giả vờ đổi nét chữ đi chứ.''

Bài tiếp theo: bài ghi tên Shouichi, thực chất là bài gốc của Joe. Cậu chàng Ma cà rồng thở dài. Bài làm này nham nhở, các phép toán rối rắm, lời giải lộn xộn như một trang giấy nháp.

``Đây đúng là phong cách của gã hề\dots\ Dù có là Shouichi viết thì cũng không tệ đến mức này.''

Cuối cùng là bài ghi tên Lize, thực chất là bài gốc của Shouichi.

Chỉ mới đọc qua một nửa, Kuzuha đã phải bóp trán. Nét chữ ẩu, viết ngoáy, không có lời giải mà chỉ toàn đáp án chép lụi, có câu thậm chí\dots\ ghi sai cả đề.

``Đúng là mấy kẻ không học bài có khác.''

Cậu Ma cà rồng rũ vai, rời khỏi đống lộn xộn để tìm chút thanh lọc tinh thần bằng cách mở hai bài mà cậu trông chờ nhất: Hikawa Kuon và Hikawa Kaigou - hai kẻ im lặng suốt giờ kiểm tra, không nói, không gian lận, không phô trương gì.

Đầu tiên là bài của Kuon. Toàn bộ phần đại số và giải tích được lý giải bằng những lập luận cực kỳ gãy gọn. Cách trình bày như thể cậu ta đang viết tài liệu hướng dẫn, không thừa một dòng nào.

Nhưng khi tới phần hình học - sở đoản của cậu - nét chữ bắt đầu ngập ngừng hơn, lời giải dài dòng hơn, kèm theo cả ghi chú như: ``Tôi không chắc lắm nhưng có vẻ đúng nếu áp dụng định lý Menelaus\footnote{Định lý phát biểu như sau: Cho tam giác ABC bất kỳ. Các điểm D, E, F lần lượt nằm trên các đường thẳng BC, CA, AB. Khi đó D, E, F thẳng hàng khi và chỉ khi $\left|\frac{\overline{FA}}{\overline{FB}}\right| \times \left|\frac{\overline{DB}}{\overline{DC}}\right| \times \left|\frac{\overline{EC}}{\overline{EA}}\right| = 1$}''. Dẫu vậy, đáp án cuối cùng\dots\ lại chính xác.

Cậu nhếch miệng cười thầm: ``Ít ra cậu ta cũng thành thật với phần yếu của mình.''

Tiếp đến là bài của Kaigou. Cô nàng nổi bật với hình học không gian, và điều đó hiện rõ rành rành trong từng dòng viết: cô tính từng thể tích, từng mặt cắt, từng góc chiếu với độ chính xác đến từng mi-li-mét. Kèm theo là\dots\ bốn tờ giấy nháp, kín đặc ký hiệu véc-tơ, mặt phẳng, và hình chiếu.

Tới phần đại số - điểm yếu của cô - những dòng giải bắt đầu trở nên rối hơn, nhiều trường hợp phải xét riêng. Thậm chí, một bảng biến thiên còn được vẽ hẳn ở ô trả lời bên cạnh, tới mức mà Kuzuha cũng không nhịn nổi cười. Nhưng, giống như Kuon, Kaigou vẫn đưa ra đáp án đúng, và thậm chí còn đính kèm biểu đồ tuyến tính được vẽ tay cực kỳ cẩn thận.

``Không hổ là cặp song kiếm tương phản,'' Kuzuha nghĩ thầm, ánh mắt lóe lên một tia thích thú. ``Chắc chắn mình sẽ tung kết quả vào mặt bọn chúng, từng đứa một.''

Cậu đóng bài lại, khoanh tay nhìn đống bài thi trước mặt mình. Trong đầu Kuzuha, màn kịch tiếp theo sắp được dựng lên, và dĩ nhiên, người dẫn dắt màn kịch ấy, không ai khác ngoài cậu chàng Ma cà rồng lắm tài nhiều tật ấy.

\ct

Một tiếng sau.

Tất cả hai mươi người được gọi quay lại lớp học để nhận kết quả bài kiểm tra. Ai nấy đều thấp thỏm, riêng Kuon và Kaigou vẫn ngồi thản nhiên, ánh mắt nhìn ra cửa sổ như thể chuyện này chẳng ảnh hưởng gì đến họ.

Ngay lúc ấy, một giọng nói quen thuộc vang lên từ cái đồng hồ đeo tay của Kuon. ``\eng{Mà này, sao sáng nay cậu không đeo tôi theo như thường lệ? Tự dưng tôi thấy yên tĩnh quá\dots}''

``\eng{Làm bài kiểm tra,}'' cậu lười biếng đáp.

``\eng{\dots Ra thế. Chẳng trách.}''

``\eng{Thì đương nhiên. Làm bài kiểm tra thì có được phép mang gì ngoài giấy bút đâu?}''

Kaigou cũng chen vào: ``\eng{Với cả nếu mà có đeo thì chắc chắn ông cũng sẽ nói là: `Đừng bao giờ nghĩ đến việc gian lận thi cử!' hay đại loại như thế\dots}''

``\eng{\dots Cũng có lý.}''

Ở phía trên bục giảng, Kuzuha đang cầm trên tay xấp bài kiểm tra. Cậu gõ nhẹ vào bàn: ``Được rồi lớp, đây là kết quả bài kiểm tra của mọi người. Tôi sẽ đọc từng người. Nếu ai nghe thấy tên mình thì đứng lên nhé.''

Cậu rút ra một tờ đầu tiên, mắt nhìn lướt qua cái tên Lize Helesta được ghi ngay ngắn bên dưới dòng tên Joe Rikiichi bị gạch đỏ. Con số màu đỏ nằm ngay trên hình vẽ con chim non ở góc bài. ``Lize-san, tám mươi lăm điểm.''

Cả lớp ồ lên. Cô nàng tóc trắng đứng bật dậy, đôi mắt long lanh như sắp khóc, miệng lẩm bẩm: ``Mình\dots\ mình làm được rồi sao? Thật không thể tin được!'' Cô quay vòng tại chỗ, váy công chúa xoè ra như đang biểu diễn, ánh sáng lấp lánh từ cửa sổ rọi vào trông như một cảnh trong anime.

\begin{figure}[H]
    \centering
    \begin{subfigure}[t]{0.45\textwidth}
        \centering
        \includegraphics[width=8cm]{../../../../Image/Book?/nj-8l1.png}
        \caption{Kết quả bài thi của Lize (được chỉnh lại tên) với 85 điểm.}
    \end{subfigure}
    \quad
    \begin{subfigure}[t]{0.45\textwidth}
        \centering
        \includegraphics[width=8cm]{../../../../Image/Book?/nj-8l2.png}
        \caption{Phản ứng của cô nàng.}
    \end{subfigure}
\end{figure}

``\vie{Cũng không phải dạng vừa đâu nhỉ,}'' cậu Tổng liếc về phía nàng Công chúa đang nhảy chân sáo lấy bài.

``\vie{Công chúa thật sự không đùa đâu,}'' cô nàng ``song sinh'' gật đầu đồng tình.

Mashiro, ngồi phía sau cô nàng, trong lòng nở hoa: ``\textit{Tám lăm à\dots? Nếu mình chép y chang thì chắc mình cũng được tám lăm\dots\ Hoặc hơn!}'' Joe và Shouichi liếc nhau, rướn mày, trong đầu đang tính toán điều gì đó mờ ám.

``Tiếp theo là Rikiichi-san và Kanda-san.''

Cậu cầm nhẹ hai tờ bài lên. Mỗi tờ đều có một con số đỏ choét in tròn góc phải, và một cái tên được sửa lại to rõ: ``Zê-rô. Không điểm.''

Âm thanh đó như hai tiếng sét ngang tai. Joe và Shouichi lập tức bật dậy, hai tay ôm lấy mặt, run rẩy.

``Không thể nào\dots\ không thể nào\dots!''

``Chẳng lẽ mình tráo nhầm bài!?''

\begin{figure}[H]
    \centering
    \begin{subfigure}[t]{0.45\textwidth}
        \centering
        \includegraphics[width=8cm]{../../../../Image/Book?/nj-8m1.png}
        \caption{Kết quả bài thi của Joe và Shouichi (được chỉnh lại tên) với 0 điểm.}
    \end{subfigure}
    \quad
    \begin{subfigure}[t]{0.45\textwidth}
        \centering
        \includegraphics[width=8cm]{../../../../Image/Book?/nj-8m2.png}
        \caption{Phản ứng của hai người khi biết kế hoạch bị vỡ lở.}
    \end{subfigure}
\end{figure}

``\vie{\dots Biết ngay là kiểu gì cũng có mấy thành phần kiểu này mà.}''

``\vie{Ừ. Còn lạ gì nữa.}''

``Bài của hai cậu chỉ có tên là khác nét chữ thôi,'' Kuzuha bình thản nói.

``Khoan đã, cái gì cơ?'' Kuon ngạc nhiên.

``Đừng nói là\dots'' Kaigou chưa kịp nói hết câu.

``\eng{Đúng như hai người nghĩ đấy. Hai cậu ấy tráo bài cho nhau,}'' Game lạnh lùng xác nhận, thở dài một tiếng.

``Không dễ qua mặt tôi đâu,'' cậu Ma cà rồng chốt lại, giọng khản đặc nhưng sắc lẹm.

``Thật là vãi chưởng\dots'' cả hai người nhà Hikawa đều cùng đồng thanh, mắt trố ra. ``Không ngờ tới mức đó luôn đấy.''

``\eng{Đây là ông thầy tinh ý nhất mà tôi từng thấy,}'' Game nhận xét.

Kuzuha lật sang bài tiếp theo: ``Tiếp theo là\dots\ Mashiro-san.'' Cậu chàng tóc đen lập tức sáng rực mắt, môi nở nụ cười tự tin. Nhưng chỉ sau hai giây, cậu bị bắn trở lại mặt đất bởi hai từ chí mạng: ``Zê-rô.''

\begin{figure}[H]
    \centering
    \begin{subfigure}[t]{0.45\textwidth}
        \centering
        \includegraphics[width=8cm]{../../../../Image/Book?/nj-8n1.png}
        \caption{Phản ứng mong chờ của Mashiro.}
    \end{subfigure}
    \quad
    \begin{subfigure}[t]{0.45\textwidth}
        \centering
        \includegraphics[width=8cm]{../../../../Image/Book?/nj-8n2.png}
        \caption{Phản ứng sau khi biết Kuzuha bắt cậu tại trận.}
    \end{subfigure}
\end{figure}

Một màu đen đột ngột trùm xuống, hai con ma trơi trắng lượn quanh đầu Mashiro. Mắt cậu trợn ngược, miệng há hốc ra: ``Không\dots\ không thể nào\dots''

``Err\dots\ Cụ thể là cậu ta đã làm gì vậy?'' Kaigou hỏi.

``Đáp án của cậu ấy hoàn toàn trùng với đáp án của Lize-san\dots'' Kuzuha đáp.

``Và\dots?''

``\dots Kể cả những đáp án sai.''

Kuon và Kaigou phì cười. ``Bảo sao.''

``0 điểm là cũng đáng.'' Game nhận xét.

Ba người quay xuống nhìn thấy Lize đang giận sôi máu, ánh mắt lạnh lẽo như mùa đông Bắc Âu. Cô đứng dậy, đi tới trước bàn của Mashiro - người đang thè lưỡi, mặt xanh lè như lá chuối - và giơ nắm đấm lên\dots

\gif{nj-8o}{1}{164}

*BỐP!*

Một cú đấm không cần lời giải thích. Cậu chàng tóc đen bị đấm bay khỏi ghế, cổ vặn vẹo dài ngoằng, thân thể lôi thẳng ra khỏi lớp như một sợi mì Ý bị quăng đi.

``\vie{Cái quái gì vậy\dots}''

``\vie{Cô ấy mạnh thật đấy!}

``\eng{Và cái cổ của cậu ta thì thật không bình thường\dots}''

Lize đứng xoay một vòng, váy công chúa vẫn tung bay, hai tay chống hông như thể vừa giành chiến thắng vinh quang. Cô không nói một lời nào, chỉ quay lưng về chỗ ngồi, mặt lạnh tanh.

\img{nj-8p}

Cậu Ma cà rồng thở dài: ``Lize-san, em về chỗ đi. Tôi biết em là người liêm khiết, nhưng\dots''

``Như thế có hơi quá không, Công chúa Helesta?'' giọng Pháp cất lên, tỏ vẻ e ngại cho cậu Mạo hiểm vừa bị hạ đo ván kia.

Lize vẫn không đáp lại. Cô lẳng lặng ngồi xuống ghế, khoanh tay đầy kiêu hãnh.

``Được rồi, ờm\dots\ Tiếp theo này.''

Kuzuha rút ra hai tờ bài cuối cùng trong xấp giấy kiểm tra, đảo mắt qua cặp đôi Hikawa đang ngồi phía trên.
``Hikawa-san, cả hai cô cậu.''

``Vâng?'' cặp đôi đồng thanh đáp lại, nửa tò mò, nửa chờ đợi.

Cậu chàng Ma cà rồng giơ cao cả hai tờ bài, để toàn bộ lớp có thể nhìn rõ con số đỏ chót được viết ở góc phải mỗi bài thi: 100.

``Cả hai người đạt điểm tuyệt đối. Cũng là hai người nghiêm túc nhất trong lớp này.''

Kuon nở nụ cười nhẹ, gật đầu một cái đầy hài lòng nhưng không quá phô trương: ``Tôi cũng đoán được phần nào rồi. Nhưng mà được trăm điểm thì đúng là\dots\ hơi sốc đấy.''

Kaigou thì trợn mắt, khẽ thốt lên: ``Ơ thật á!? Tôi\dots\ tôi được trọn điểm á!?'' Cô như không tin vào tai mình, nhưng rồi liền vươn người đập tay với Kuon. ``Lần đầu tiên tôi đạt điểm tối đa trong bài kiểm tra toán đấy!''

``Tốt hơn là cô nên cảm thấy tự hào,'' Kuon cười trừ.

Kuzuha lúc này ngồi xuống mép bàn giáo viên, đặt xấp bài còn lại xuống và rút riêng ra một tờ trong đó. ``Trong bài kiểm tra này, có một câu tôi cố ý gài bẫy. Đó là một câu hỏi dạng phân số trong phần giải hệ phương trình.''

Cậu giơ lên đề thi, rồi chỉ vào phần câu hỏi mà giờ đây ai cũng nhận ra là ``có vấn đề'': dấu gạch phân số in mờ. ``Chính là chỗ này. Tôi cố tình làm cho dấu gạch phân số này mờ đi một chút.''

Cả lóp xôn xao một hồi. Cậu chàng tiếp tục: ``Tôi muốn kiểm tra khả năng phân tích thị giác và logic. Dù chỉ là một chi tiết nhỏ, nhưng nếu các em đủ tinh ý, thì vẫn có thể suy luận ra đâu là phân số, đâu là phép trừ.''

Cậu đưa mắt nhìn sang cặp đôi ngồi đầu bàn. ``Cả hai người đã nhận ra được điểm bất ổn này. Giải thích cho mọi người xem nào.''

Cậu Tổng đứng lên trình bày: ``Tôi thấy phần tử số được căn giữa, khoảng cách phía dưới có dư, nên tôi loại trừ khả năng đây là dấu trừ. Hơn nữa, dấu đó quá dài để được gọi là phép trừ, và dấu trừ thường được đặt ở trên tử, và nếu là phép trừ, dấu gạch ấy phải dài hơn một chút sang trái.''

Cô nàng ngồi cạnh cậu thì bật cười: ``Còn tôi thì\dots\ nhầm đấy. Ban đầu tôi tưởng là dấu trừ thật, nên làm thử. Nhưng khi ra kết quả thì nó lẻ quá mức. Mà tôi nhớ đề toàn câu chọn đáp án đẹp cơ mà? Thế là tôi thấy không ổn, làm lại theo kiểu phân số\dots\ và may là kết quả ra đúng.''

``Và đó là điều khiến hai người vượt qua cái bẫy mà rất nhiều người khác mắc phải. Bao gồm\dots'' Kuzuha nhìn về phía bàn ở sát cửa sổ, ``\dots Lize-san.''

Cả lớp quay lại nhìn. Công chúa tóc trắng ngồi im, hai mắt mở to không thể tin nổi khi bị người giám thị kia chỉ điểm. Cô sực nhớ ra cô đã làm như vậy mà không hề kiểm tra lại bài.

``Không thể nào\dots'' cô khẽ nói, giọng gần như vỡ ra vì bàng hoàng. ``Làm sao mình lại sai một thứ đơn giản thế được!? Mình\dots\ mình là người được huấn luyện quân sự chính quy mà!''

Cô đưa tay ôm trán, miệng thì thào như thể đang tự chất vấn bản thân mình vì một sai sót ngớ ngẩn đến không tưởng.

Kuon nghiêng người sang thì thầm với Kaigou: ``\vie{\hl{Tôi tưởng công chúa thì phải bắn súng giỏi, hóa ra bị hạ bởi một dấu phân số mờ\dots}}''

Kaigou suýt bật cười nhưng cố giữ vẻ nghiêm túc: ``\vie{\hl{Cẩn thận không cô ấy nghe được đấy\dots}}''

Kuzuha không bình luận gì thêm, chỉ lặng lẽ cầm bài tiếp theo, nhưng trong ánh mắt của cậu, rõ ràng là rất thích thú với kết quả từ hai ``kẻ liêm khiết'' kia, những người không chỉ tránh được mọi chiêu trò, mà còn vượt qua được chính cái bẫy nhỏ mà cậu dày công sắp đặt.

%==========================================TRUYỆN PHỤ=========================================
\omk

Tiết cuối cùng kết thúc.

Lize vẫn ngồi lặng người ở bàn, ánh mắt mơ hồ nhìn vào khoảng không trước mặt như chưa hoàn hồn sau cú sốc. Mãi đến khi chuông reo báo hiệu kết thúc buổi học, cô mới từ từ đứng dậy. Vẫn còn chút ngập ngừng, cô cầm lại tờ bài kiểm tra của mình rồi bước nhanh về phía cửa lớp, nơi Kuon và Kaigou đang chuẩn bị rời đi.

``Khoan đã\dots\ Hikawa-san.''

Giọng nói đầy cương quyết vang lên từ sau lưng, khiến cặp đôi nhà Hikawa quay lại. Lize Helesta - công chúa của một vương quốc, và giờ là thí sinh vừa sảy chân trong một cái bẫy đơn giản - đang đứng trước họ, hai tay chống hông, đôi má đỏ bừng vì xấu hổ hơn là giận dữ.

``Tôi không phục lắm đâu đấy\dots'' cô nàng khẽ hắng giọng.

Cậu Tổng nhướng mày, nửa cười nửa tò mò. ``Cô không kiểm tra lại bài à?''

``Có chứ!'' công chúa phồng má, quay mặt đi một chút. ``Đến mức có đủ thời gian để vẽ cả Sebastian đấy thôi!''

Cô rút tờ bài kiểm tra ra, đưa lên trước mặt họ. Trên góc phải, một hình vẽ khá đáng yêu của con chim non hiện rõ - đôi mắt sáng tròn, chiếc nơ đỏ nhỏ xíu dưới cổ, và cái mỏ hếch lên đầy kiêu hãnh. ``Sebastian,'' cô nói với vẻ nghiêm trang, ``là quản gia trung thành của tôi.''

``Cô có vẻ yêu mến lắm nhỉ?'' Kuon hỏi với sự tò mò, lần này là sự chân thành, không móc mỉa như lúc nãy.

``Đúng là như vậy\dots''

Nghe đến đây, giọng của Game vang lên với chất mỉa mai đặc trưng: ``Tôi nghĩ\dots\ thay vì phí thời gian vẽ con chim kia thì có thể dành nó để kiểm tra lại bài đấy, công chúa.''

``\textbf{Ông im đi!}'' Lize bặm môi đáp lại, quay ngoắt đi chỗ khác như muốn tránh cả thế giới.

Kaigou cười khúc khích, vỗ vai cô một cái nhẹ nhàng. ``Thôi nào. Tám mươi lăm điểm vẫn là tốt rồi. Trong lớp này, nhiều người còn không qua nổi trung bình nữa là.''

Cô liếc nhìn Mashiro đang nằm sõng soài dưới đất, một tay ôm mặt, tay còn lại gãi đầu trong vô vọng, như thể vẫn chưa tin mình vừa ăn một cú đấm từ chính người bị sao chép.

Ở bàn gần đó, Joe và Shoichi thì gục đầu xuống bàn như hai cái cây chuối đổ rạp sau bão.

Gã Hề thì thở dốc như vừa chạy marathon, thì thầm bằng chất giọng lạc đi: ``Sao mình lại tin thằng kia chứ\dots'' còn cậu Truyền thông thì đập đầu xuống bàn cái ``cộp'', lẩm bẩm: ``Mình nên học hành tử tế thì hơn\dots''

Lize liếc một vòng ba kẻ gian lận ấy, rồi quay sang Kaigou: ``Ừ, chắc là cậu nói đúng.'' Cô hít một hơi sâu, ngẩng đầu nhìn thẳng về phía trước, giọng trở lại như thường lệ: ``Lần sau, tôi sẽ lấy lại thể diện.''

Kuon nheo mắt cười. ``Thì lần sau cứ đừng vẽ chim nữa là được.''

``Cậu-!''

Một lần nữa, má của nàng Công chúa lại ửng đỏ, không rõ vì giận hay vì xấu hổ. Nhưng lần này, thay vì phản pháo, cô lặng lẽ quay lưng đi, giữ cho cái kiêu hãnh của vương tộc không vỡ vụn trước mặt hai anh em nhà Hikawa.

Kaigou nghiêng đầu nhìn cô, nở một nụ cười thân thiện: ``Này, nếu cậu rảnh\dots\ qua ký túc xá bọn tôi chơi không? Tiện thể để Kuon sửa lại bài cho. Còn mấy câu khó, tôi cũng có thể giúp một tay.''

Lize thoáng ngạc nhiên. Có vẻ cô vẫn chưa quen với kiểu mời mọc thẳng thắn như vậy, nhất là khi vừa mới lãnh một cú sốc tâm lý: cả từ bài thi lẫn cú đấm trước đó. Nhưng cô cũng nhanh chóng lấy lại bình tĩnh, chỉnh lại mái tóc dài óng ánh và đáp lời bằng giọng hoàng gia đặc trưng:
``Được thôi. Dù gì thì\dots\ tôi cũng không định thua mãi đâu.''

``Cách tốt nhất để đánh bại được đối thủ là phải học được cách đối thủ đánh nhau như thế nào, đúng chứ?'' Kuon hỏi.

``Ừ, cậu nói chí phải.''

\ct

Thế là nhóm ba người rời khỏi lớp học, về lại khu ký túc. Suzuki, lúc này đang lau bàn uống nước ở phòng khách, ngẩng đầu lên khi thấy hai anh chị trở về.

``Onii-sama! Onee-sama! Làm bài thi xong rồi ạ?''

``Xong rồi,'' Kuon gật đầu, rồi rút trong túi ra tờ đề thi của mình. ``Em thử nhìn câu này xem, thấy có gì sai sai không?''

Cậu chỉ tay vào chỗ khoanh dấu đỏ, chính là phần phân số bị làm mờ trong đề. Cô bé cầm đề lên, nghiêng đầu một chút rồi nhíu mày: ``Ơ\dots\ Đây rõ ràng là dấu phân số mà? Chứ sao lại là dấu trừ được?''

Kuon cười nhẹ: ``Ừ, anh cũng thấy thế.''

*CHOANG!*

Lize đứng bất động. Đôi mắt nàng công chúa nhìn chằm chằm vào Suzuki, như thể vừa bị ai đó lấy búa đập một cú chí mạng vào lòng tự tôn.

Cô bé này\dots\ thậm chí không học cùng cấp bậc, không được huấn luyện nghiêm ngặt như cô, không có quân sư riêng, không có thư viện riêng\dots\ mà lại nhìn ra được thứ mà cô - một người từng qua hàng trăm bài kiểm tra học thuật của hoàng tộc - lại không thể nhận ra.

``Sebastian\dots'' nàng Công chúa khẽ nói như lẩm bẩm với bản thân – ``\dots có lẽ ta nên bớt vẽ cậu đi một chút\dots''

Game  lại không bỏ qua cơ hội: ``Công nhận em gái Kaigou-user nhìn chuẩn thật. Nếu để em ấy chấm bài, có khi cũng chẳng ai chối cãi được.''

Suzuki ngẩng đầu: ``Em chấm kiểu cộng thẳng điểm à?''

``Không, ý là\dots\ Thôi, bỏ đi.''

Lize vẫn chưa nói gì thêm. Cô chỉ ngồi im, đặt đề thi xuống bàn, rồi bất chợt ngẩng lên nhìn về phía cửa sổ. Mặt trời dần lặn sau những dãy nhà nhiều màu sắc của vùng đất Cầu Vồng, để lại ánh sáng cam nhạt phủ khắp phòng.

Dường như cả bốn người trong phòng - Kuon, Kaigou, Suzuki và cả giọng nói trong chiếc đồng hồ của cậu Tổng - đều nghĩ đến cùng một điều.

Một ngày nữa lại trôi qua.

Và họ vẫn chưa thể quay về.

``Ngày thứ tám rồi đấy,'' Kaigou nói nhỏ, tựa lưng vào ghế.

``Cũng đã hơn một tuần trôi qua rồi, Onee-sama\dots,'' Suzuki đếm bằng ngón tay, ánh mắt mơ màng.

``Chỉ mong có lỗ xoáy nào đó tới kéo chúng ta đi luôn,'' Game càu nhàu.

Kuon không nói gì. Cậu chỉ chống cằm nhìn về phía xa, giọng trầm hẳn xuống: ``\dots Rời khỏi nơi này càng sớm càng tốt.''

\end{document}